%%%%%%%%%%%%%%%%%%%%%%%%%%%%%%%%%%%%%%%%%%%%%%%%%%%%%%%%%%%%%%%%%%%%%%
% Lớp 8 | Chương X. Một số hình khối trong thực tiễn | Hình chóp tứ giác đều
% Dạng 8H0?39-1: Nhận biết các yếu tố của hình chóp tứ giác đều
%%%%%%%%%%%%%%%%%%%%%%%%%%%%%%%%%%%%%%%%%%%%%%%%%%%%%%%%%%%%%%%%%%%%%%


% nguồn: De_AI_Toan8_cuoi_hoc_ki_i_1 (nhập 21/11/2026)
\begin{bt}%[8H0H39-1]%[ma:ecf26f27]
	Cho hình chóp tứ giác đều $S.ABCD$ có cạnh đáy bằng $18\text{ cm}$ và chiều cao $SO = 12\text{ cm}$. Gọi $M$ là trung điểm của cạnh $BC$.
a) Tính độ dài các đoạn thẳng $OM$ và $SM$.
b) Tính diện tích xung quanh của hình chóp $S.ABCD$.
	\loigiai{
		a) Vì $O$ là trung điểm của $AC$ và $M$ là trung điểm của $BC$ nên $OM$ là đường trung bình của tam giác $ABC$, suy ra $OM = \dfrac{AB}{2} = \dfrac{18}{2} = 9\text{ cm}$. \\
Áp dụng định lí Pythagore trong tam giác $SOM$ vuông tại $O$: $SM = \sqrt{SO^2 + OM^2} = \sqrt{12^2 + 9^2} = \sqrt{225} = 15\text{ cm}$. \\
b) Diện tích xung quanh của hình chóp là $S_{xq} = (2 \cdot 18) \cdot 15 = 36 \cdot 15 = 540\text{ cm}^2$.
	}
\end{bt}

% nguồn: De_AI_Toan8_cuoi_hoc_ki_i_1 (nhập 21/11/2026)
\begin{bt}%[8H0H39-1]%[ma:66a26b9b]
	Cho hình chóp tứ giác đều $S.ABCD$ có cạnh đáy bằng $14\text{ cm}$ và chiều cao bằng $24\text{ cm}$.
a) Tính độ dài trung đoạn của hình chóp.
b) Tính diện tích toàn phần và thể tích của hình chóp.
	\loigiai{
		a) Khoảng cách từ tâm đáy đến trung điểm một cạnh đáy là $r = \dfrac{14}{2} = 7\text{ cm}$. \\
Độ dài trung đoạn là $d = \sqrt{24^2 + 7^2} = \sqrt{576 + 49} = \sqrt{625} = 25\text{ cm}$. \\
b) Diện tích xung quanh: $S_{xq} = (2 \cdot 14) \cdot 25 = 700\text{ cm}^2$. \\
Diện tích mặt đáy: $S_{\text{đáy}} = 14^2 = 196\text{ cm}^2 \Rightarrow S_{tp} = 700 + 196 = 896\text{ cm}^2$. \\
Thể tích hình chóp: $V = \dfrac{1}{3} \cdot 196 \cdot 24 = 1568\text{ cm}^3$.
	}
\end{bt}
